\documentclass[11pt]{article}
\usepackage[margin=1in]{geometry}
\usepackage{fontspec}
\setmainfont{Noto Serif CJK SC}
\setsansfont{Noto Sans CJK SC}
\setmonofont{Latin Modern Mono}
\XeTeXlinebreaklocale "zh"
\XeTeXlinebreakskip = 0pt plus 1pt minus 0.1pt
\usepackage{amsmath,amssymb}
\usepackage{booktabs}
\usepackage{array}
\usepackage{graphicx}
\usepackage[dvipsnames]{xcolor}
\usepackage[colorlinks=true,linkcolor=NavyBlue,citecolor=PineGreen,urlcolor=NavyBlue]{hyperref}
\usepackage{xspace}
%% 由 scripts/make_tables.py 生成，勿手改；取值来自 results/c5_matrix.json 与 results/r37e_full_seeds.json。
%% 正文与表说明里的这些数字只允许写成这些宏。
\newcommand{\cFiveAllSparseYellowA}{547}
\newcommand{\cFiveEnergyDeadA}{0}
\newcommand{\cFiveEnergyDeadB}{0}
\newcommand{\cFiveEnergyLostA}{290}
\newcommand{\cFiveEnergyLostB}{826}
\newcommand{\cFiveEnergyYellowA}{559}
\newcommand{\cFiveEnergyYellowB}{1194}
\newcommand{\cFiveNightfloorDeadA}{20}
\newcommand{\cFiveNightfloorDeadB}{31}
\newcommand{\cFivePeak}{0.012}
\newcommand{\cFiveSeeds}{3}
\newcommand{\cFiveTtlEightDeadA}{0}
\newcommand{\cFiveTtlEightDeadB}{0}
\newcommand{\cFiveTtlEightRevertA}{10}
\newcommand{\cFiveTtlEightRevertB}{9}
\newcommand{\cFiveTtlEightYellowA}{849}
\newcommand{\cFiveTtlEightYellowB}{2020}
\newcommand{\cFiveTtlFourLostB}{104}
\newcommand{\cFiveTtlFourRevertA}{6}
\newcommand{\cFiveTtlFourRevertB}{5}
\newcommand{\cFiveTtlFourYellowA}{849}
\newcommand{\cFiveTtlFourYellowB}{1916}
\newcommand{\cFiveValidUntilYellowA}{849}
\newcommand{\cFiveValidUntilYellowB}{2020}
\newcommand{\cFiveYellowNA}{2016}
\newcommand{\cFiveYellowNB}{3528}
\newcommand{\cThreeDeathsFifo}{12}
\newcommand{\cThreeDeathsPurge}{0}
\newcommand{\cThreeExpiredFifo}{2538}
\newcommand{\cThreeExpiredPurge}{120}
\newcommand{\cThreeLatestDelta}{5.55}
\newcommand{\cThreePairedHi}{4.60}
\newcommand{\cThreePairedLo}{3.48}
\newcommand{\cThreePairedPoints}{4.04}
\newcommand{\cThreeRatio}{2.45}
\newcommand{\cThreeSeedZeroFifo}{202}
\newcommand{\cThreeSeedZeroPurge}{420}
\newcommand{\cThreeSvcFifo}{0.373}
\newcommand{\cThreeSvcPurge}{0.413}
\newcommand{\cThreeTenSeedFifo}{1691}
\newcommand{\cThreeTenSeedPurge}{4145}
   % 由 scripts/make_tables.py 生成，勿手改
\linespread{1.12}
\tolerance=2400
\emergencystretch=3.2em
\hbadness=10000
\vbadness=10000
\newcommand{\stime}{$S_{\mathrm{time}}$\xspace}
\newcommand{\scap}{$S_{\mathrm{cap}}$\xspace}
\newcommand{\senergy}{$S_{\mathrm{energy}}$\xspace}
\newcommand{\saccess}{$S_{\mathrm{access}}$\xspace}
\newcommand{\unk}{$U$\xspace}
\newcommand{\trel}{t_{\mathrm{rel}}}

\title{\textbf{间歇回传下灾前监测的本地通信控制}}
\author{灾前监测与应急通信工作稿\\ \texttt{adaptive-communication} 仓库 —— 工作稿}
\date{}

\begin{document}
\maketitle

\begin{abstract}
间歇回传会切断监测系统的控制路径，却不会停止此前安装的通信状态：未确认记录继续占用接入机会，已生效的采样配置在纠正命令无法到达后继续耗电。本文研究电池与光伏供电、LoRaWAN Class~A 接入、蜂窝主回传与仅上行短报文备用组成的灾前山区监测网，区分记录生命期、授权配置回退与端到端完成证据，并研究各类处置能够执行的位置。设计依据是：有效的局部行动所需信息可以少于完整根因诊断。在所测缓存模型中，标准期限到期安置在源端后，十种子按期服务配对提高 \cThreePairedPoints 个百分点，中断按期交付达到基线的 \cThreeRatio 倍；仅在网关抑制发送则可能反压接入。本地保护也能避免所测场景中中心降档无法及时阻止的损失。普通 TTL 与电量规则覆盖了受检配置停止问题，包括同信息精确参照，因此这些结果不依赖新租约算法。网关局部归因与真实 Agent 轨迹进一步暴露了把接入收据理解为回传或完成证据的接口故障。现有证据支持本地执行与明确证据边界的系统设计原则；Agent 研究仍属形成性验证，完整执行接口相对同组件普通组合的独立增益尚待检验。
\end{abstract}

\noindent\textbf{关键词：}灾前监测；间歇回传；LoRaWAN Class~A；通信状态；本地执行；期限到期；部分可观测；Agent 执行。

\section{引言}
山区监测节点在回传正常时收到一条授权密采配置。蜂窝路径随后中断，节点继续采样和上报，后续降档命令却无法到达 Class~A 接收窗口。与此同时，未获得中心确认的记录留在源缓存，即使已过业务期限仍占用后续接入批次。网关停止发送无价值记录，也未必能让源端副本释放。某一层的正确决定因此可能无法阻止另一层的持续资源消耗。

本文研究\textbf{控制失联下持续执行的通信状态}。记录保留与已生效配置各有终止条件，却共享一个执行依赖：等待中心后续处置需要一条可能已经中断的路径。监测义务仍是端到端的——窗口内的合格样本必须在期限内到达中心；证据与执行权限却分散在节点、网关和中心。接入收据只能确认其中一段。

\subsection{设计依据}
\textbf{局部处置与完整诊断具有不同的信息要求。} 网关可能无法区分从未采集与尚未听到的样本，源端仍能确定某条记录是否超过本地已知期限并停止重传。授权回退也可以依据本地时钟或电量动作，即使中心尚不清楚当前档位。处置需要明确的权限和可执行的本地条件，无须臆测远端故障原因。

本文将这种区分组织为执行契约：标识动作影响的状态、支持结果判断的证据，以及来得及干预的位置。标准到期和本地保护是契约的具体实现。系统价值通过安置位置及通信后果评估，不由成熟机制的新名称决定。

\subsection{普通机制与剩余问题}
对照包含标准逐记录到期、本地时钟与电量保护、固定有效期回退、强备份装包以及同信息精确停止参照。标准到期复现了所测种子中按义务索引的清理行为；普通配置规则覆盖了受检停止族。这些结果限定了新算法的必要性，但不构成所有通信调度器的全局最优性证明。

剩余系统问题是：面对已经拥有这些普通组件的通信栈，显式执行接口能否进一步改善动作选择与修复。Agent 可能在缺口位于回传时继续增加采样，也可能用无关接入收据推断命令已生效。真实轨迹已观察到后一类接口故障，其端到端比较尚未隔离接口的独立贡献。该问题与确定性的执行位置效果分别验证。

\subsection{贡献}
\begin{enumerate}
\item \textbf{面向失联监测的执行契约}（\S\ref{sec:model}--\S\ref{sec:mech}）。区分本地可判定的记录到期、授权配置回退与部分端到端证据，以证据可得、权限和动作时延共同描述执行位置条件。
\item \textbf{跨段安置效果的实证}（\S\ref{sec:eval-record}、\S\ref{sec:eval-place}）。源端标准到期带来 \cThreePairedPoints 个百分点的服务增益（95\% 区间 [\cThreePairedLo,\cThreePairedHi]），中断按期交付达到 \cThreeRatio 倍；配置对照说明同一保护规则在中心无法执行、在节点仍可动作的条件。
\item \textbf{控制与证据的经验边界}（\S\ref{sec:eval-lease}--\S\ref{sec:eval-agent}）。同信息普通停止规则追平声明的精确参照，网关在线归因只能覆盖部分义务，形成性 Agent 轨迹暴露接入与回传证据混淆。这些对照区分了已有本地补救和待验证的规划器或运行时增量。
\end{enumerate}
研究始终限定于外部授权任务下的灾前监测；系统不自行判断预警等级，不删除失约义务，也不通过更改评分任务取得额外通信资源。

\section{场景与工作点}
\label{sec:bg}

十四个传感器节点分两组监测坡体位移与降雨。节点经 LoRaWAN Class~A 上传到现场网关，网关经蜂窝回传转发到中心服务器；北斗短报文作为稀疏、计量的备用链路，只承载现场到中心的消息，不充当下行配置通道。节点由电池加小光伏供电，任务是灾前例行监测并带授权的预警等级变更，任何组件都不判滑坡风险。节点数、角色映射与扰动参数按冻结的 Task v1.1 实例清单，早期的 16 节点草图不作为现场事实。

授权边界来自两项中国地质灾害监测行业标准 DZ/T~0460--2023 与 DZ/T~0450--2023，条款在仓库内逐条审计~\cite{dzt0460,dzt0450}。\S5.3.3 允许按预警等级远程动态调整采样与上传频率且未给数值，因此本文使用的周期是已声明的研究选择；\S8.4.1 定义红橙黄蓝四级，\S8.4.2 允许会商后升级、降级或解除；\S5.3.7 在高风险极端天气弱覆盖下要求双模通信，该条只出现一次，属条件性条款，通信层标准并未通篇强制。北斗公开服务规范 BDS-OS-PS-3.0~\cite{bdsps3} 给出成功率 $\geq$95\%、平均时延小于 2~s、平均每 30~s 一次服务、单条消息上限 14{,}000 比特，同时说明实际频度与最大长度受注册参数约束，频度属于配额量。这些条款划出合法动作范围：在授权等级内重新编译频率并申请会商变更，静默删除义务与自行宣告风险都在界外。

备用链路每 1200~s 一个 78 字节包（20~B 包头、58~B 载荷；位移样本 6~B、降雨 4~B，每包约装 9--10 个）。升级后的黄级任务周期 300~s，交付期限为窗口开启后两个周期，即交付窗约 600~s，短于 1200~s 的备份槽间距。DZ/T~0450 的 200~B 单包上限与不小于 1~min 的周期（\S7.4.3.2、\S7.4.2.3，\S7.4.3.5 另有不小于 5~s 的消息间隔）分属不同的服务形式条款，不合并成单一的"北斗工作点"；1200~s/78~B 是已声明的研究选择。

\section{系统模型}
\label{sec:model}

时间以 $\Delta{=}60$~s 为槽，$t{=}0$ 为本地 06:00，故 $t{=}12$~h 是日落，日照周期在模型里显式出现。节点集 $\mathcal N$、测量量集 $\mathcal M$ 与网关 $g$ 构成现场，中心 $c$ 在回传之外。

\subsection{义务与交付链}
监测义务 $o=(n,m,[l_o,h_o),d_o)$ 定义在节点 $n$、测量量 $m$ 上，含采样窗口与绝对交付期限。周期为 $P$ 的例行义务有 $h_o-l_o=P$、$d_o=h_o+P$，即一个周期的宽限。样本满足 $o$ 的条件是测量量匹配、在窗口内采集、并在 $d_o$ 之前被中心收到，依次经过
\begin{equation}
\resizebox{\linewidth}{!}{$\displaystyle
\underbrace{\text{合格样本}}_{\text{节点}}
 \!\to\! \underbrace{\text{在 }g\text{ 被听到}}_{\text{LoRa 接入}}
 \!\to\! \underbrace{\text{在 }d_o\text{ 前回传}}_{\text{蜂窝/短报文}}
 \!\to\! \underbrace{\text{中心证据}}_{\text{义务满足}} .$}
\end{equation}
分段任务的首段用闭区间、后续段用半开窗口，等级变更后边界样本不会重复覆盖相邻窗口。

\subsection{链路、中断与可行动实体}
LoRa 接入（节点到网关，含 Class~A 下行窗口）每次机会以 $p_a{=}0.74$ 成功；主回传（网关到中心）在非声明中断期以 $p_b{=}0.62$ 良好。\textbf{控制中断}是区间 $\mathcal O=[a,b)$，其间主回传对中心到现场的流量不可用：中心命令与任务表段在准入处被拒、从未发出，短报文备用无法承载下行配置。中断期间能改变现场通信行为的只有\textbf{节点}（依据自身时钟、电量、缓存、最近接入收据，以及已生效命令携带的界）与\textbf{网关}（依据实际听到的样本、转发日志与已用备份槽）。中心保留授权变更的权力，没有安装变更的通道；在 $\mathcal O$ 内撰写的修订只能在 $b$ 时刻到达现场。

\subsection{持久安装状态}
\textbf{记录状态。}节点 $n$ 上的每个未确认样本 $s$ 有生成时刻 $t_s$，在本地已知节奏下业务期限 $d(s)=(\lfloor t_s/P\rfloor+2)P$，占用 $K{=}240$ 个缓存槽之一，在一次接入批次中最多占 $B{=}32$ 个位置。在所建模的边缘纪律下，它只在中心确认时离开缓存（确认按同步、不占空口建模，对各队列规则同等有利），在此之前按最老优先重发。

\textbf{配置状态。}节点对采样与上报运行同一个档 $q(t)\in\{P_0,P_1\}$（稀疏 $P_0{=}600$~s、密集 $P_1{=}300$~s），由带版本号的命令安装。节点收到并执行命令时配置生效，后续上报可向观察者确认该效果；排队或途中的命令不改变配置。生效期间 $q{=}P_1$ 每拍额外消耗采样与上行能量，直到触发终止。

\subsection{能量}
单次采样耗 $e_s{=}4.7{\times}10^{-4}$~Wh，单次上行耗 $e_u{=}2.33{\times}10^{-5}$~Wh，电池容量 $C{=}0.05$~Wh。日照输入是 12~h 半正弦，峰值 $\eta$（标称 $0.03$~Wh/h），叠加逐节点逐小时的随机云遮挡。主配置下电量归零的节点永久死亡。整夜维持密集采样消耗约 $0.071$~Wh，超过电池容量。

\subsection{本地判据与授权回退}
记录携带按源端已知节奏赋予的到期时刻 $d_n(s)$。当 $t>d_n(s)$ 时，它对该本地契约所表示的义务已无交付价值；截止当拍仍可参与传输。任务修订尚未抵达时，节点不能把未来中心期限作为本地输入，独立评分仍保留实际授权义务。

配置则具有预先允许的回退集合及本地触发条件，例如已安装 TTL、时钟规则或电量条件。停止新增密采与删除已采记录是两个动作：配置改变后，已有样本仍可能有用。因此本文不再用最后回传槽定义统一的配置释放时刻；未知的未来任务修订不能成为节点停止的证据。

\subsection{服务损失与声明纪律}
目标函数计三项：按期交付的义务（服务）、节点永久死亡、以及存活超过释放时刻的状态所占资源——浪费的样本、接入尝试、短报文字节与能量。离线口径下，每条未满足义务记在乐观模型（后续假设一律放行）下的首个失败环节上，得到四个互斥段：\senergy（无合格样本）、\saccess（采到但期限前未被网关听到）、\stime（已听到但不存在回传机会）、\scap（存在机会但容量或拥塞承载不了）。在线观察者看不到某一环节时报告 unknown，不得猜测。每条可行性声明属于三态之一：\textbf{confirmed}（独立上报显示目标档，或有匹配样本被中心收到）、\textbf{unconfirmed}（命令在途中或收据不全，既不蕴含已生效也不蕴含链路已断）、\textbf{unreachable}（由明确的网络侧状态得知当前路径已断）。LoRa 接入收据只构成接入链路的证据，凭它本身永远不能确立回传可达。

\section{问题形式化}
\label{sec:problem}

令 $\mathcal H_e(t)$ 只包含实体 $e$ 在 $t$ 前实际持有的证据。节点历史包括时钟、电量、缓存和已收到命令；网关还持有听到的样本与自身转发日志。策略通过 $a_e(t)=\pi_e(\mathcal H_e(t))$ 决定记录保留或释放、配置维持或回退。比较结果向量为
\begin{equation}
\mathcal J(\pi)=\bigl(N_{\rm miss},N_{\rm brownout},E_{\rm node},A_{\rm up},A_{\rm down},B_{\rm backup}\bigr),
\label{eq:obj}
\end{equation}
分别报告失约、掉电、节点能量、上下行空口和备用字节。单节点参照中的死亡权重是声明的实验偏好，不是从灾害监测任务推出的常数。全部未来存活是可选风险口径，并非每个实验必须满足的隐含条件。

\subsection{执行契约}
对照遵守四项条件。\textbf{局部性}：相同本地历史不能仅因不可见未来不同而产生不同动作。\textbf{授权}：回退动作必须位于事先授予的权限内。\textbf{保留责任}：源端可按已知生命期释放记录，但释放传输状态不能从评分中删除义务。\textbf{结果证据}：请求排队、配置生效、网关听到和中心交付分别建立不同事实。

契约保留实际取舍。回退可以保护能源，同时减少未完成义务的采集；损失仍计入 $N_{\rm miss}$。仅更换控制器无法保证所有资源轨迹下同时满服务与存活。

\subsection{处置放在哪里}
对判据 $c$ 与执行层 $e$，设计条件为
\begin{equation}
\operatorname{evidence}(e,c)\land\operatorname{authority}(e,c)
\land D_{\rm observe+decide+act}(e,c)<M_c(t),
\label{eq:placement}
\end{equation}
其中 $M_c(t)$ 是剩余干预余量。有证据、有权限仍不足以使中心保护规则可执行：其时延包含回传与 Class~A 窗口，可能超过能量余量。源端到期无需新命令或远端故障诊断。该条件说明实现要求；经验时延估计不构成普遍安全界。

\subsection{可知边界}
考虑节点观测完全相同的两段历史：一段中心已降级但消息未到，另一段高等级要求仍有效。在新证据到达前，任何因果节点策略在两段历史中采取同样动作，因此无法声称精确跟踪不可见修订。预授权回退规定节点可以做什么，却不揭示远端究竟发生了哪段历史。这个不可区分性说明未知结果需要与已确认完成或不可行分开。

Agent、确定性规则和求解器都可使用同一契约；共同动作空间与证据边界定义了独立于规划器的通信问题。

\section{机制}
\label{sec:mech}
原型在三处实例化契约：源端执行已知记录生命期与授权回退，网关使用实际听到的样本和有限回传机会，中心编译要求并解释收到的证据。下列实验分别隔离组件；完整组合是否优于拥有相同机制的普通通信栈仍需独立比较。

\subsection{记录终止：按期限对齐的到期}
\label{sec:mech-record}
每次上传机会上，节点删除 $d_n(s)<t$ 的缓存记录，其余按最老优先发送。实现两种写法：\texttt{deadline-purge} 以义务节奏为索引，\texttt{generic-expiry} 是普通 BPv7 式逐记录生命期，在生成时赋值 $\text{expires}=t_s+((P-t_s\bmod P)+P)=d(s)$，不引用义务 id、槽几何或证书。网关变体在本地确认它将不发送的记录。EDF 与义务贪心队列只重排、从不释放；AoI 的 \texttt{latest-only} 无差别释放，只保留最新一条。丢弃陈旧或过期流量在 AoI 包管理与 DTN bundle 到期中已有充分积累~\cite{costa2014pm,kam2016deadline,dtnarch,rfc9171}，抑制超出时延界的传感流量同属先行工作~\cite{wada2009timeout}。本文不主张新的丢弃算法。系统层的内容是端到端业务期限基准（用业务期限，而非包龄、逐跳 TTL 或估计的单包时延）、零下行的源端安置，以及决定性的跨段一致性——确认耦合使得仅在网关释放并不能释放源端缓存。

\subsection{普通保护下的配置回退}
\label{sec:mech-lease}
本地时钟门夜间把采样与上报钳到稀疏，白天保留配置。配置对照另评估固定 TTL、滚动电量与期望采能检查、能量推导停止界。TTL 从实际空口密采生效边沿起算；各臂共享源端到期、备份装包和允许的回退档位。

区分两种信息条件。预告有效期场景中，节点随配置获得已知终点；运行中授权场景中，中断期间才发布的降级在现场不可见，本地规则使用时钟、实测电量和声明的能量模型。预告有效期具有额外信息，不能作为后者的同信息竞争者。另用公开任务表下的量化单节点停止模型，将普通规则与精确非预知参照比较，范围和结果见\S\ref{sec:eval-lease}。

\subsection{部分可观测投影器}
\label{sec:mech-projector}
编译期中心持有义务表、当前（从不预测）的网络侧中断状态与公开参数；任务表到达时网关只持有它听到过的东西。结构性无回传证书（\stime）仅由槽几何产生，并显式以主回传不可用为条件，几何本身从不断言中断。到达时刻投影器只在本地位证据能区分各段时给出标注，否则返回 unknown。unknown 感知的中心投影器只在存在明确网络侧字段时输出 unreachable（当前 harness 尚未提供该字段，故输出 unconfirmed），并用算出的能量台账替代一律夜间关停，使真正的能量不可行与保守建议区分开来。

\section{原型与方法}
\label{sec:proto}

全部机制运行在同一个 Python 事件步进仿真器里（60~s 一拍），实现 LoRaWAN Class~A 下行窗口、带版本的双字段配置命令（采样间隔与上报周期作为独立下发字段，跨版本混合配置会被检出）、确认前保留的边缘缓存、最大覆盖备份装包、随机日照与云遮挡、吸收态掉电，以及非预知的分段任务视图门。关闭备份可逐位复现 v1.1 仿真器（五个锚定不变量）；本地门与回退扩展默认关闭。

使用两个受控场景。\textbf{Episode~I} 是 v1.1 清单工作点：$t{=}6$~h 升级到黄级，主回传中断 4--20~h，任务表在 20~h 到达网关，驱动前沿、归因、记录到期、执行位置与真实 agent 结果。\textbf{Episode~II} 是确定性的配置终止场景，除时刻表、中断相位与被扫描的采能外参数相同：为把白天降档放进中断之内，相位 A 取 $t{=}2$~h 升黄级、$t{=}6$~h 会商降级、中断 4--20~h；相位 B 取 $t{=}1$~h 升黄级、$t{=}8$~h 降级、中断 6--22~h；$t{=}12$~h 为日落。Episode~II 是为配置子问题显式构造的受控对照，不构成新的现场主张。节点死亡按种子 0--2 的总数报告；$n{=}3$ 下亚百分点级的平均服务差按不显著处理；种子 6--9 是重复使用的留出确认集，不是未被触碰的测试集。表格由在册结果文件生成，复现入口提供确定性检查和冻结参考。

\section{评估}
\label{sec:eval}

\subsection{固定资源下的损失位置}
\label{sec:eval-frontier}
Episode~I 的 7560 条例行义务上，强 maxcov/FIFO 基线交付 3025（0.400）。对称的逐义务反事实每次只放宽一项物理资源、从不改策略，在相同义务 id 上比较集合差。回传放宽到无限得 4740（0.627），能量放宽（0.50~Wh、全程密集）得 4089（0.541），两者同时放宽得 6023（0.797）；仅被回传容量阻塞 1715 条，仅被电池阻塞 1128 条，两者共同阻塞 155 条（2.0\%），见表~\ref{tab:walls}。这些是在已实现策略上的干预读数，不是穷尽最优；本文不作零残差或全局最优主张。它们定位了该工作点的资源压力；状态终止实验处理同资源下的额外执行后果，其动机不依赖证明所有调度器都已耗尽优化空间。

\begin{table}[t]
\centering\small
\caption{7560 条义务上的反事实干预（Episode~I）。在已测配置上的集合分解，不是全部合法策略的上界。}
\label{tab:walls}
% 由 scripts/make_tables.py 生成，勿手改；数字来自结果文件。
\begin{tabular}{llcc}
\toprule
配置 & 采样/能量 & 回传 & 交付\\
\midrule
base & dayfeed, 0.05 Wh & 1200 s/78 B & 3025 (0.400)\\
$O_{bh}$ & dayfeed, 0.05 Wh & 无限 & 4740 (0.627)\\
$O_{en}$ & always-300, 0.50 Wh & 1200 s/78 B & 4089 (0.541)\\
$O_{bo}$ & always-300, 0.50 Wh & 无限 & 6023 (0.797)\\
\bottomrule
\end{tabular}

\end{table}

把 4535 条失约按首个失败环节记账，且"听到"必须发生在期限之前（而非此后任意时刻），得到 \senergy~2002 条（44.1\%）、\scap~1115 条（24.6\%）、\saccess~830 条（18.3\%）、\stime~588 条（13.0\%）。此前的"曾经听到"口径报 \scap~1945、\saccess~0：被挪动的 830 条实际上都采到了，只是首次被网关听到晚于期限，属于接入迟到，与回传容量无关。两种口径下交付数（3025）相同；纠正改变的是损失记在哪一段，因而改变哪条机制能够处理它。

\subsection{记录到期：标准原语，位置决定成败}
\label{sec:eval-record}
Episode~I 中断期间 75 个备份包承载 750 条记录；FIFO 下 219 条到达中心时已经过期，按期 261 条（45.6\%；2782 载荷字节中用掉 1270）。仅在网关抑制会去掉这些发送，却扣住确认，按期交付从 202 掉到 194，代价被搬到接入段。源端按期限到期改变符号：中断按期交付 \cThreeSeedZeroFifo 提到 \cThreeSeedZeroPurge。十种子下（表~\ref{tab:purgeexpiry}）到期把平均服务从 $\cThreeSvcFifo$ 提到 $\cThreeSvcPurge$，配对 $+\cThreePairedPoints$ 点（95\% 区间 $+\cThreePairedLo$ 到 $+\cThreePairedHi$，十个种子全部为正）；中断按期交付总数 \cThreeTenSeedFifo 升到 \cThreeTenSeedPurge（\cThreeRatio 倍），过期备份记录 \cThreeExpiredFifo 降到 \cThreeExpiredPurge，死亡 \cThreeDeathsFifo 降到 \cThreeDeathsPurge；相对 AoI \texttt{latest-only} 领先 $+\cThreeLatestDelta$ 点，因为 latest-only 用放弃恢复后覆盖换取中断期新鲜度。\texttt{deadline-purge} 与 \texttt{generic-expiry} 在每种子上逐位相同（服务、中断计数、死亡与交付集合全部一致）：记录侧机制就是标准的逐记录到期，它的实测价值来自业务期限基准与源端网关跨段一致性，与新的丢弃算法无关。

\begin{table}[t]
\centering\small
\caption{固定强 maxcov 选择器的十种子 Episode~I 扫描，只变边缘处置。仅网关抑制（未列，同 FIFO 缓存）在种子 0 得 194，对比 FIFO 的 \cThreeSeedZeroFifo 与源端到期的 \cThreeSeedZeroPurge，即跨段代价。$^{\dagger}$标准 BPv7 式逐记录生命期；$^{\ddagger}$按义务索引的核销，与 $^{\dagger}$ 逐位相同。}
\label{tab:purgeexpiry}
\setlength{\tabcolsep}{4pt}
% 由 scripts/make_tables.py 生成，勿手改；数字来自结果文件。
\begin{tabular}{lcccc}
\toprule
边缘处置 & 全时段服务 & 中断按期 & 过期备份 & 死亡\\
\midrule
FIFO（基线） & 0.373 & 1691 & 2538 & 12\\
latest-only（AoI） & 0.358 & 4250 & 0 & 0\\
\texttt{generic-expiry}$^{\dagger}$ & 0.413 & 4145 & 120 & 0\\
\texttt{deadline-purge}$^{\ddagger}$ & 0.413 & 4145 & 120 & 0\\
\bottomrule
\end{tabular}

\end{table}

\subsection{配置保护与算法边界}
\label{sec:eval-lease}
Episode~II 固定源端到期与 maxcov 装包，对比不同本地回退。时钟门在日落动作，固定有效期和能量规则可提前动作。预告有效期具有提前任务终点信息，与中断期间无法获知降级的策略分列。

\begin{table*}[t]
\centering\small
\caption{两个中断相位的配置回退对照（\cFiveSeeds 个种子，固定 \texttt{generic-expiry} 与 maxcov）。TTL 从密采生效起算。预告有效期行为任务~1 信息条件，其余自适应与 TTL 行为任务~2。三条已报告能量变体逐格结果相同，故只列一行。实测零死亡不构成普遍存活保证。}
\label{tab:lease}
% 由 scripts/make_tables.py 生成，勿手改；数字来自结果文件。
\begin{tabular}{lcccc}
\toprule
本地退回界 & 死亡（3 种子） & 平均服务 & 最终 SoC & 黄级交付\\
\midrule
\multicolumn{5}{l}{\emph{相位 A，阴雨 $\eta{=}0.012$（升级 2~h、降级 6~h、中断 4--20~h）}}\\
时钟门（日落 12~h） & 20 & 0.5961 & 0.066 & 849\\
固定 TTL 4~h（退回 6~h） & 0 & 0.5970 & 0.122 & 849\\
固定 TTL 6~h（退回 8~h） & 0 & 0.5970 & 0.122 & 849\\
固定 TTL 8~h（退回 10~h） & 0 & 0.5971 & 0.117 & 849\\
\textbf{交付租约 ($\tau_L{=}6.33$~h)} & \textbf{0} & \textbf{0.5970} & \textbf{0.122} & \textbf{849}\\
\midrule
\multicolumn{5}{l}{\emph{相位 B，阴雨 $\eta{=}0.012$（升级 1~h、降级 8~h、中断 6--22~h）}}\\
时钟门（日落 12~h） & 31 & 0.5981 & 0.038 & 2020\\
固定 TTL 4~h（退回 5~h） & 0 & 0.5881 & 0.122 & 1837\\
固定 TTL 6~h（退回 7~h） & 0 & 0.5994 & 0.122 & 1994\\
固定 TTL 8~h（退回 9~h） & 0 & 0.5996 & 0.122 & 2020\\
\textbf{交付租约 ($\tau_L{=}8.33$~h)} & \textbf{0} & \textbf{0.5996} & \textbf{0.122} & \textbf{2020}\\
\bottomrule
\end{tabular}

\end{table*}

$\eta=\cFivePeak$ 时，时钟门在 A、B 相位留下 \cFiveNightfloorDeadA 与 \cFiveNightfloorDeadB 个死亡。固定 8~h TTL 为零死亡，黄级交付为 $\cFiveTtlEightYellowA/\cFiveYellowNA$ 与 $\cFiveTtlEightYellowB/\cFiveYellowNB$；4~h TTL 在 B 相位少交付 \cFiveTtlFourLostB 条。受检能量门同样零死亡，但相对 8~h TTL 分别少交付 \cFiveEnergyLostA 与 \cFiveEnergyLostB 条。其预测负载相对实测射频负载偏保守，这组结果限定于该实现，不能否定全部预测能量控制。矩阵支持本地回退的作用，未提供新能量推导租约的独立增量。

\textbf{同信息停止参照。} 为区分实现局限与优化空间，量化单节点模型在完整物理地平线上选择继续密采或吸收式退回稀疏。公开任务表决定奖励，云遮未来仍未知；模型按当前仿真采样与射频负载标定，普通 TTL、电量阈值与滚动规则获得相同信息。在 \cNineCells 个相位、采能与离散度格及 \cNineWeights 个死亡权重上，普通组合在 \cNineServiceMatches/\cNineComparisons 项比较中与精确非预知参照服务一致，计价目标在数值精度内相同。声明的紧能量点上，全部未来存活的严格端点对全程稀疏也不可行；这是风险模型边界，不表示监测任务或所有物理实现不可行。

该负结果关闭受检单节点停止族，不构成新的配置算法，也不确立共享网络调度或未知任务表下的最优性。普通回退继续作为执行系统组件。

\subsection{任务表到达时的可判定范围}
\label{sec:eval-attr}
Episode~I 中延迟的任务表在 $t_a{=}20$~h 到达网关，此时 2136 条黄级义务已过期限。离线时间感知归因把 2136 条全部标对（\stime~588、\scap~327、\saccess~650、\senergy~571）。只用网关在 $t_a$ 之前合法持有的证据，可标注 1138/2136（53.3\%）——全部几何 \stime（588 条）、全部按期听到但未回传的 \scap（327 条）、以及 223 条可证迟到的 \saccess——在已标注子集上精度 100\%；998 条（46.7\%）保持 unknown，其离线真值是 571 条从未采样、427 条在 $t_a$ 之后才被听到，见表~\ref{tab:attribution}。"从未采样"与"采到但尚未听到"在缺少网关手里没有的证据时无法分离，必须保持 unknown，这正是原始的 G5 约束。编译期几何不需要样本，输出 602 条条件性无回传证书，其中 588 条为真 \stime，14 条恢复边界情形不予认证、直接排除；乐观能量容量证书不纯（552 条正确 \scap、571 条实为 \senergy、193 条假阳），因此容量判定留给网关。已知期限下的源端到期可以在远端原因未明时执行；配置回退则另需事先授权。

\begin{table}[t]
\centering\small
\caption{任务表到达网关时已经错过 2136 条升级义务的到达时刻归因（Episode~I）。在线投影器精确但部分；427 条匹配样本在 $t_a$ 之后才首次被听到，650 条晚于各自期限。}
\label{tab:attribution}
\setlength{\tabcolsep}{4pt}
% 由 scripts/make_tables.py 生成，勿手改；数字来自结果文件。
\begin{tabular}{lcc}
\toprule
段 & 离线真值 & 在线（$t_a$）\\
\midrule
\stime（无回传槽，几何） & 588 & 588\\
\scap（按期听到但未回传） & 327 & 327\\
\saccess（晚于期限听到） & 650 & 223\\
\senergy（无合格样本） & 571 & --- (unknown)\\
\unk（在线不可分离） & --- & 998\\
\midrule
合计 & 2136 & 2136\\
标注覆盖率 & 100\% & 53.3\%\\
已标注子集精度 & 1.000 & 1.000\\
\bottomrule
\end{tabular}

\end{table}

\subsection{执行位置决定存活}
\label{sec:eval-place}
中心侧可行性信封（对任何规划器做纯过滤，输入只用已确认配置、听到计数、授权要求与时钟）去掉命令风暴，却救不了最紧的种子：日落降级命令必须穿过回传与一个它错过的 Class~A 窗口。同一条"勿整夜密集"规则在中心是迟到、可被拒的建议，在节点成为每拍可执行的本地钳制（表~\ref{tab:r39}）。朴素中心合规在三个种子上都杀死 14 个节点中的 13 个（合计 39 个），服务 0.162；加上节点本地门后全部死亡消除而服务不变；有门在场时，信封进一步把被拒准入尝试压到 222。零命令预置节奏达到 0.360 服务，与这些静态任务上各中心编译臂相当——这是针对预置任务的中心采样控制的限定负结果；普通本地控制器应进入后续 runtime 与 Agent 比较的共同底座。

\begin{table}[t]
\centering\small
\caption{确定性 Episode~I 执行位置研究（无 LLM）。本地门在所测种子上消除全部永久死亡且服务不变，仅中心信封做不到。被拒命令在准入处退回、从未发出，属命令卫生而非计费节省。点特定经验结果，不是生存不变量或最优界。}
\label{tab:r39}
\setlength{\tabcolsep}{4pt}
% 由 scripts/make_tables.py 生成，勿手改；数字来自结果文件。
\begin{tabular}{lcccc}
\toprule
配置（除注明外均在中心） & 平均服务 & 死亡 & 发送 & 被拒\\
\midrule
comply（朴素，无保护） & 0.162 & 39 & 160 & 2488\\
dayfeed-c（中心昼夜规则） & 0.370 & 12 & 308 & 1256\\
envelope(comply)，仅中心 & 0.359 & 12 & 308 & 298\\
comply + 本地门 & 0.366 & 0 & 1424 & 3364\\
dayfeed-c + 本地门 & 0.368 & 0 & 258 & 1156\\
\textbf{envelope(comply) + 本地门} & \textbf{0.359} & \textbf{0} & \textbf{282} & \textbf{222}\\
\midrule
\textbf{pure-local（预置，0 命令）} & \textbf{0.360} & \textbf{0} & \textbf{0} & \textbf{0}\\
\bottomrule
\end{tabular}

\end{table}

\subsection{形成性真实 agent 研究}
\label{sec:eval-agent}
在 Episode~I 中断上实跑 deepseek-flash：无辅助 agent（A0）、带普通待决命令清单的 agent（A0s）、带证书的 agent（A1）各三个种子，另加两个健康链路对照；每 30 分钟一次真实模型调用，每条轨迹 99 次，合计 1089 次决策。各臂共享 harness、合法观测、工具、历史与动作空间。这项研究用于激发并压力测试接口，不用来确立确定性机制。

\textbf{服务与存活（$n{=}3$，噪声大）。}A1 平均服务 0.358，与 0.360 的确定性安全底线持平，死亡 4；A0 为 0.296、死亡 35；A0s 为 0.222、死亡 39；健康链路对照下 A0 为 0.275、死亡 13，A1 为 0.511、死亡 0。逐种子服务区间重叠，故不就服务作显著性判断。实跑 agent 回路未接本地门，因此 A1 在种子 1 的 4 个死亡不由 \S\ref{sec:eval-place} 的确定性结果清除。

\textbf{调用账。}1089 次决策中有 52 次因 \texttt{max\_tokens} 截断 JSON 而解析失败、退回 hold。逐臂解析失败率为 A0 7.1\%（28/396）、A0s 6.4\%（19/297）、A1 1.3\%（5/396），因此臂间差异混合了输出格式可靠性与语义效应。

\textbf{接口故障与对称评分。}观测的 \texttt{link} 字段只承载 LoRa 接入收据（\texttt{n\_heard}），不含回传状态。强制中断窗内，经区分词边界的声明检查，A0/A0s 分别在 36/32 次决策上断言「link 在线、命令可以生效」，A1 为 1 次。主回传全程中断，因此这些声明过度乐观，与无效密采命令和死亡同时出现。反方向也存在：v4 证书在听到节点数不足全员时断言「控制面未交付」，在主回传实际可用的五个健康链路时刻触发，其中两次发生在 13/14 个节点已经确认密集档之后。声明在中断与健康条件下均受检查，不以该决策是否下发命令为前提。

对相同快照离线重放 unknown 感知的 v5 投影器，缺少网络侧字段时不作不可达断言，五个健康链路假警报全部变为 unconfirmed；528 次一律夜间关停被拆成 189 次能量不可行与 339 次建议性，临近黎明的 04:30--06:00 全部属于建议性。缺少回传状态信号会造成双向错误。隔离投影器的独立增量仍需同能力专家规则对照、带显式回传状态字段的端到端 v5 运行及跨模型复现。

\section{相关工作}
\label{sec:related}

\textbf{应急与网络控制中的 LLM。}TopoLLM 把自适应工具学习用于灾后拓扑规划~\cite{topollm}；后续系统让 LLM 调度无人机路径与带宽、塑造 RL 奖励或编排恢复工具，涉及意图对象、预算与时限约束、候选方案检查、局部修复以及有文档的网络管理 API~\cite{icgrestore,ossgpt}。自然语言转约束与工具编排已有这些基础。本文关注持续现场状态、Class~A 生效时延和独立中断的回传路径，引用上述工作以定位控制器接口；未运行其原始系统，不据此认领复现或否定它们。

\textbf{调度、新鲜度与能量。}状态模型上最近的先例是带每传感器有界队列、外生到达、有限服务机会与反馈的 LLM 上下文内数据采集调度器~\cite{icldc}；AoI 与能量采集控制，含瞬时电池阈值策略与部分观测采集，处理的是资源片段~\cite{bacinoglu1905,zakeri2311,costa2014pm}；空间链路的超额预订与接触感知转发~\cite{ccsdssabr} 与网络演算界~\cite{leboudec01} 是互补的确定性刻画。联合的采样、传输与能量优化已经确立，本文不作主张；本文的区分在于带 Class~A 生效时延的持久远端配置、分离的接入与主备回传路径，以及控制中断下的终止。一致网络更新研究早已处理逐更新的有效性、依赖关系与安全迁移~\cite{sdnconsistency}；本文沿用这一教训，补充断开且能量受限的边缘情形：已安装更新的自身数据在反向更新到达之前就变得不可交付。

\textbf{包丢弃、DTN 到期与接触。}AoI 管理抢占或丢弃旧更新~\cite{kam2016deadline,krishnan2019multihop}；面向时延的缓冲丢弃已与自适应 ARQ 结合用于视频~\cite{razavivideo}；DTN bundle 协议携带生成时刻的生命期并按 TTL 或相遇丢弃~\cite{dtnarch,iranmaneshdtnqueue,rfc9171}，托管传输~\cite{rfc5050} 把保留责任作为前向交接转移（该机制已从 BPv7 核心迁出）。乡村 LoRa 与低轨卫星系统围绕估计的卫星接触排程并利用网关空间分集~\cite{sateriot}。本文的记录处置就是 BPv7 式到期并给出等价性证明；配置侧评估普通本地回退。已支持的系统内容是带源端网关跨段一致性的端到端业务期限基准，以及提前本地回退相对日落回退的实测效果。同信息停止参照已被普通组合覆盖。托管转移是中心 ACK 保留纪律的实质替代；改变反馈与存储责任时，它应继续作为强对照。

\section{讨论与局限}
\label{sec:disc}

\textbf{已建立的系统认识。} 源端到期、本地保护与部分归因各有独立证据，共同指向处置位置应保有行动所需的证据与能力。最强正结果来自源端保留责任，配置停止则说明普通机制足够的范围。未来更强的调度器不会消除已经观察到的缓存滞留和降档不可达现象。

\textbf{组合与迁移。} 完整执行 runtime 需要与具有相同 expiry、本地保护、影子、观测和规划工具的普通组合比较，不能把组件效果相加作为独立增益。现有结果覆盖声明部署与受控片段，中断相位、接入与回传共同故障、采能条件决定迁移边界。所需本地固件、期限元数据、信令与存储代价以及硬件执行尚未由仿真验证。

\textbf{反馈与保留责任。} 当前中心确认同步且不占空口，同等有利于各队列规则。仅上行备用链路下，实际中心交付未必立即成为源端知识。部署需要区分接入收据、网关持久托管与中心完成，释放通知如果经过 Class~A 窗口就应付出真实成本。标准托管转移、累计确认与源端到期因此都是必要普通对照。有限反馈与配置容量的联合分配尚未进入本文评估。

\textbf{任务与能源范围。} Episode~I 源端按本地已知的变更前节奏赋予期限，中心仍按实际任务评分。多义务共享记录时可能需要保留到最后一次有效使用。采能是合成过程，主配置把掉电建模为吸收态。精确停止参照具有公开任务表和量化能量模型，不能将其最优性外推到未知任务或完整通信网。

\textbf{Agent 证据与剩余接口问题。} 实跑为单模型，格式可靠性与专家信息包存在混杂；已支持的结论是接口失效及投影器的离线修正。同能力端到端比较应共享本地保护、证据、确定性工具与解析预算，检验动作接口能否减少有害或无效执行，同时保留有益变更，并独立核对完成声明。跨模型复现应跟随已经隔离的作用；增加调用量本身不能建立该作用。

\textbf{来源。} 标准条款按仓库获取件交叉核对，一条复制链在官方端点不可达时使用文档分享站。部署参数区分已公开能力与研究选择。本文未执行被引系统原实现，不主张复现其性能。

\section{结论}
间歇回传会移除控制能力，而已经安装的通信状态仍继续消耗资源。在所测灾前监测网中，标准源端到期提高按期服务 \cThreePairedPoints 个百分点，中断交付达到基线的 \cThreeRatio 倍；本地保护能在后续中心降档不可达时继续动作。普通规则覆盖了受检配置停止问题。由此形成的设计原则是：以本地可得证据、授权和可执行位置组织保留与回退，同时将端到端完成与中间收据分开。Agent 可以使用这种执行接口，其独立增益仍需同能力普通组合对照。现有证据支持通信机制及其适用边界，不把规划器智能程度作为正确本地行动的前提。

\bibliographystyle{IEEEtran}
\bibliography{../refs}

\end{document}
